\subsection{The direct unit-trace criterion in the recovered source}
\label{sec:step17-unit-interface}

This section uses a stronger intermediate statement actually proved in
step~17's source, rather than trying to recover it from the eventual
existential estimate. All source declarations below are at the pinned
commit of \cite{K17PinnedLeanMirror}.

\begin{lemma}[K17-S17-UNIT]\label{lem:step17-unit}
Assume the source datum $\mathrm{GD}(\Omega,K,\ell,e,\psi)$.
Write $\Psi=\chi_\ell^e\psi$, interpreted in the coefficient DVR, and
fix $s\in G_K$ and $t=\Psi(s)+1+\Psi(s)^{-1}$.
If, for every source-admissible representation
\[
 \rho\in\mathrm{Adm3U\_s}(\iota,e,N)
       \ \cup\ \mathrm{Adm2A}(\iota,e,N|d_K|),
\]
the difference $\operatorname{tr}\rho(s)-t$ is an $\ell$-adic unit,
then the source Selmer length $V_{\mathrm{len}}(\Omega,c,e,\psi)$ is zero.
\end{lemma}

\begin{proof}
The exact statement \texttt{OneCharU\_s} is defined in
\texttt{DenseThresholdCases/NumberU3US.lean}.
It requires only nonvanishing of all the indicated differences and concludes
that $V_{\mathrm{len}}$ is finite and that a finite set $X$ of those
admissible representations exists with
\[
 \frac{V_{\mathrm{len}}}{e(O/\mathbf Z_\ell)}
 \le 2\sum_{\rho\in X}
       v_\ell\bigl(\operatorname{tr}\rho(s)-t\bigr).
\]
The normalization is the definition \texttt{Chr.nrm} in
\texttt{CharacterLevelBounds/ShaFamMono.lean}.
The statement is proved by
\texttt{oneCharU\_of\_katoTwoW\_rig\_admI\_automorphy} in
\texttt{DenseThresholdCases/FOneCharU.lean}, including its preceding
complete proof \texttt{oneChar\_unramU}. The proof constructs the actual
finite set from the model's \texttt{TwoPt} points, treats Type~I and
Type~II with the two displayed predicates, and removes the unramified
coefficient restriction by \texttt{exists\_unram} and
\texttt{coefIndep}.

The same closed source inputs supply this intermediate theorem:
\texttt{SmallClassNumbers/CorGKOf.lean}, in the proof of
\texttt{corGK\_of}, applies precisely that declaration with all its
cohomology, algebra, Kato, rigidity, admissibility and automorphy arguments.
Its subsequent \texttt{litN} and \texttt{corGK} supply the closed published
inputs. We retain this intermediate theorem value, rather than passing it
through the congruence-number construction. This uses the stipulated source
theorem, without claiming a new Lean verification of the repository.

Units are nonzero, so the theorem applies. Every valuation in its finite
sum is zero. The ramification index is finite and positive, and the finite
module length is nonnegative. The displayed inequality therefore forces
that length to be zero.
\end{proof}

\begin{lemma}[K17-S17-U3-LEVEL]\label{lem:step17-u3-level}
Put $E=\mathbf Q(\sqrt{-3})$, $e=6$, and $N=2$.
For every odd prime $\ell$ and every representation admitted by the
source predicate $\mathrm{Adm3U\_s}(\iota,6,2)$, the associated cuspidal
$\mathrm{GL}_3(\mathbf A_E)$ representation satisfies the hypotheses of
Theorem~\ref{thm:u3-coverage}. At a split prime $p\ne\ell$, with arithmetic
Frobenius $\sigma$, its unitary Satake roots $B$ satisfy
\[
 \iota\operatorname{tr}\rho(\sigma^{-1})=\sum_{\beta\in B}\beta.
\]
\end{lemma}

\begin{proof}
We retain the source's strengthened predicates. By
\texttt{NumberU3US.lean} and \texttt{AfStmt.lean},
$\mathrm{Adm3U\_s}$ contains irreducibility, the literal level condition
$\mathrm{Level}(\ell,2,(6,0,-6),\rho)$,
$\mathrm{IsAutCsdA}(\iota,r)$ for $r=\rho\chi_\ell^{-1}$,
and the additional condition $\mathrm{Th.Unl}$.
By \texttt{IsAutCsdA.lean},
\[
 \mathrm{IsAutCsdA}(\iota,r)=
 \mathrm{AutCsd0}(\iota,r)\ \wedge\
 (r\text{ continuous}\Rightarrow\mathrm{WitA}(\iota,r)).
\]
The definition \texttt{AutCsd0} in the Foundations file
\texttt{AdmissibleResidues/IsAutomorphicOverImagQuad.lean}
is exactly a regular algebraic, conjugate-self-dual cusp representation
$\Pi$ together with its attachment to $r$.

The definition \texttt{KillsMu} in
\texttt{AutomorphicForms/BlochKatoSelmerVanishes.lean}
says that its character is trivial on automorphisms fixing all $M$-th
roots of unity. Every characteristic-zero field automorphism fixes
$\{1,-1\}$, so every \texttt{KillsMu 2} character is trivial and hence
satisfies \texttt{KillsMu 1}. Reuse the same away-prime and above-prime
level witnesses, including the same flags and local disjunctions, to get
\[
 \mathrm{Level}(\ell,1,(6,0,-6),\rho).
\]
The proved \texttt{UpperTheory.level\_twist\_cyclotomic\_pow}
in \texttt{BuildingsResidueChars/LevelTwistCyclotomicPow.lean}
then gives
\[
 \mathrm{Level}(\ell,1,(5,-1,-7),r).
\]
In particular $r$ is continuous, so its actual \texttt{WitA} hypothesis
is available.

The theorem \texttt{unlAll\_of} in
\texttt{AutomorphyLifting/AdmissibleLevelRestriction/UnlAllOf.lean}
turns the retained \texttt{Th.Unl} hypothesis into
$\mathrm{Unl1}(\rho,w)$ at each $w\mid\ell$.
By its definition in \texttt{AutomorphyLifting/WeakOfAtt.lean}, this is a
flag with diagonal characters $\kappa_i$ satisfying
$\kappa_i\ne\kappa_j\chi_\ell$ for all $i,j$.
Multiplying all the characters by $\chi_\ell^{-1}$ transports the flag
and preserves these inequalities. Thus $\mathrm{Unl1}(r,w)$ holds as well.

Apply the complete source theorem \texttt{Keep.levelCore\_alt}, or its
CSD specialization \texttt{Keep.levelKA\_alt}, in
\texttt{DenseThresholdCases/KeepLevel.lean}, with modulus one and exponent
vector $(5,-1,-7)$. The required \texttt{WitA} and above-$\ell$
\texttt{Unl1} hypotheses have both just been supplied. The CSD
specialization supplies its polarization. It concludes
\[
 \Pi_f^{K\mathrm{Lev}(E,3,1)}\ne0,
 \qquad \operatorname{InfChar}_\tau(\Pi)
       =\operatorname{infOf}(5,-1,-7)=(6,0,-6),
\]
and exact Frobenius--Satake compatibility at every $w\nmid\ell$.
The definition of $K\mathrm{Lev}$, or the proved formula
\texttt{mem\_KLev} in \texttt{AutomorphicForms/KLevTop.lean}, identifies
modulus one with $\prod_w\mathrm{GL}_3(O_{E_w})$. Therefore $\Pi$ is
unramified at every finite place, including those above $\ell$.
This deduction uses the source's level theorem with its full hypotheses;
it does not infer zero monodromy from the level disjunction alone.

In the source convention,
$\operatorname{infOf}(m)=\{(n-1)/2+m_i\}_i$.
The unitary regular cohomological infinity parameter is consequently
$\operatorname{diag}((z/\bar z)^6,1,(z/\bar z)^{-6})$, and subtracting
$(1,0,-1)$ gives compact highest weight $(5,0,-5)$.
Equivalently, $r$ has Hodge--Tate weights $(-5,1,7)$ and the source relation
is $\mathrm{HT}=1-\operatorname{InfChar}$. This checks both shifts.

Finally the source uses $\chi_\ell(\sigma)=p$ and describes Frobenius
roots of $r(\sigma^{-1})$.
The twist gives $r(\sigma^{-1})=p\rho(\sigma^{-1})$, while
\texttt{algSat 3 p} multiplies each unitary root by $p$.
Cancelling these factors proves the displayed trace identity.
\end{proof}

\paragraph{How the finite characteristic polynomial is used.}
Suppose the independent compact-space certificate gives
$\sum\beta=A_p/p^6$, and the separately established rank-two coverage
places every $\mathrm{Adm2A}(\iota,6,6)$ trace in
$1+a_p/p^6$ for an eigenform in
$S_{13}(\Gamma_0(3),\chi_{-3})$. Put
\[
 Q_p(X)=\det(XI-T_p),\qquad E_p=1+p^6+p^{12},\qquad
 N_p=(E_p-A_p)Q_p(1+p^{12}).
\]
The two trace differences at $s=\sigma^{-1}$ are respectively
\[
 \frac{A_p-E_p}{p^6},\qquad
 \frac{a_p-(1+p^{12})}{p^6}.
\]
If $\ell\nmid pN_p$, they are all $\ell$-adic units. Indeed every
characteristic root is integral, and an integral product with unit value
has unit factors at every prime above $\ell$. Repeated eigenvalues cause
no change in this argument. Lemma~\ref{lem:step17-unit} therefore applies.
This uses pointwise coverage only: it neither replaces the source's
\texttt{TrData} nor asserts that an older, larger characteristic product
divides $N_p$.
